% Full independent audits: root and descent_groups; finite certificate in logs/C_finish_curvature_length.json.
\begin{lemma}[Proved: a primary curvature-length invariant]
\label{C:curvature_length}
Let $A$ be any DVR with uniformizer $\pi$, let $F=A/\pi$,
and let $m\ge1$,
and put $R=A/\pi^m$. Fix a monic irreducible polynomial
$q\in F[z]$ with $q\ne z$.
There is an additive invariant
\[
 \nu_{q,m}(P)\in\mathbb Z/m
\]
of two-periodic complexes $P$ of finite free $R$-modules, with
respect to degreewise split short exact sequences. It is invariant
under chain homotopy equivalence and is zero if $P$ lifts to a
two-periodic complex of finite free $A$-modules.

Explicitly, lift the two differentials to $A$ and form the
curvatures
\[
 C^0=\widetilde d^1\widetilde d^0/\pi^m,
 \qquad C^1=\widetilde d^0\widetilde d^1/\pi^m.
\]
On the summands whose curvature has residual primary polynomial
$q$, these differentials are injective and their cokernels are
annihilated by $\pi^m$. The invariant is
\[
 \nu_{q,m}(P)=
   \operatorname{length}_A
     \operatorname{coker}(\widetilde d^0|_q)\pmod m.
\]
In particular the invariant applies to the two-periodic Tate
complex of any finite free $R[C]$-module, for any finite cyclic
group $C$, and gives a homomorphism
$R_R(C)\to\mathbb Z/m$ vanishing on representations lifted from $A$.
\end{lemma}
\begin{proof}
We may replace $A$ by its completion: this does not change $R$,
the residue field, or the length of any finite $R$-module.
The curvature construction is the one-equation cohomology operator
of Eisenbud \cite[Proposition~1.3 and Corollary~1.4, pp.~40--41]
{C-Eisenbud1980}. We give the needed periodic construction and all
length arguments explicitly.

Choose free $A$-modules lifting the two terms and arbitrary lifted
matrices $\widetilde d^i$. Since the original differentials compose
to zero modulo $\pi^m$, the two displayed curvature matrices exist
over $A$. Associativity and cancellation of the nonzerodivisor
$\pi^m$ give
\[
 C^{i+1}\widetilde d^i=\widetilde d^i C^i
 \quad(i\bmod2).
\]
Take the product of the characteristic polynomials of $C^0,C^1$.
Henselian factorization separates its $q$-power residual factor
from its coprime complement; see \cite[Lemma~10.153.3(3),(4),
Tag~04GG]{C-StacksHensel}. The associated polynomial idempotents
decompose both free modules, compatibly with both differentials.
Call the $q$-summands $\widetilde P_q^i$. They are finite free over
the local ring $A$. Their curvature matrices are invertible,
since $q(0)\ne0$.

On these summands the products of the differentials are $\pi^m$
times invertible matrices. Thus both differentials are injective,
and after tensoring with the fraction field they are inverse up
to invertible factors. Their source and target ranks are equal.
Moreover $\pi^m\widetilde P_q^1$ is in the image of
$\widetilde d^0$, using
$\widetilde d^0\widetilde d^1=\pi^mC^1$ and $(C^1)^{-1}$.
Consequently the indicated cokernel has finite length, is killed
by $\pi^m$, and agrees with the cokernel of $d^0$ on the corresponding
complex over $R$. This also shows that the $q$-complex over $R$ is
acyclic: if $\widetilde d^0v=\pi^mw$, then injectivity of
$\widetilde d^0$ gives
$v=\widetilde d^1(C^1)^{-1}w$, and the other degree is identical.

We next verify independence of every choice. If lifted
differentials are changed by $\pi^mh^i$, their residual curvatures
change by
\[
 C'^i-C^i=d^{i+1}h^i+h^{i+1}d^i
          \quad\text{over }R.
\]
More generally, for a chain map $f:P\to Q$, lift its matrices and
write
$\widetilde d_Q^i\widetilde f^i-
 \widetilde f^{i+1}\widetilde d_P^i=\pi^mh^i$.
Multiplying out and dividing by $\pi^m$ gives
\[
 C_Q^if^i-f^iC_P^i=d_Q^{i+1}h^i+h^{i+1}d_P^i
          \quad\text{over }R.                         \tag{C.1.28}
\]
Thus curvature is natural in the homotopy category and independent
there of the lifts and bases.

Its primary decomposition is also intrinsic in that category.
Indeed, choose any monic lift $Q(z)$ of $q$ to $R[z]$.
On a $q$-summand $Q(C)$ is nilpotent: a power is divisible by $\pi$,
and $\pi^m=0$. On a summand belonging to a different residual
irreducible it is invertible. Equation~\textup{(C.1.28)} implies
that a chain map commutes with $Q(C)$ up to homotopy. Its component
from one of these two types of summands to the other is therefore
null-homotopic, by taking a sufficiently large power of $Q(C)$
and using its inverse on the other side. The identity between two
choices consequently induces mutually inverse homotopy
equivalences of their $q$-complexes. The same argument treats
arbitrary chain homotopy equivalences.

A homotopy equivalence of two such complexes induces an
isomorphism of their $d^0$-cokernels modulo maps factoring through
finite free $R$-modules. To see this directly, apply a chain
homotopy to the degree-one terms: its part through $d^0$ vanishes
in the cokernel, and its other part factors through the free
degree-zero term. These stable isomorphisms preserve length
modulo $m$. Here is an elementary verification. If
$f:X\to Y$, $g:Y\to X$ are inverse modulo maps through free
modules and $gf-1=ab$ through a free module $E$, then
$(f,b):X\to Y\oplus E$ has left inverse $(g,-a)$.
The projection to its complementary summand $Z$ is
\[
 I-(f,b)(g,-a)=
 \begin{pmatrix}1-fg&fa\\-bg&1+ba\end{pmatrix}.
\]
Every block factors through a finite free module, using the other
relation $fg-1$ and the free summand $E$. Restriction to $Z$
shows that its identity factors through a finite free module.
Hence $Z$ is projective, and therefore free
over the local ring $R$. From
$Y\oplus E\simeq X\oplus Z$ we obtain
$\operatorname{length}_A Y\equiv
 \operatorname{length}_A X\pmod m$.
This proves independence and homotopy invariance of $\nu_{q,m}$.

Finally consider a degreewise split exact sequence of complexes.
Use the splitting to lift its middle differentials as upper block
matrices over $A$, with the chosen endpoint lifts on their
diagonals. Its curvatures are upper block matrices with those
endpoint curvatures. Apply common polynomial primary projectors
to this diagram. It gives an exact sequence of free $A$-modules
in each degree, stable under the lifted differentials. All three
$q$-differentials in degree zero are injective, so the snake lemma
gives a short exact sequence of their cokernels. Their lengths
add over $\mathbb Z$, and in particular modulo $m$.
If the original complex lifts with square-zero differential over
$A$, choose that lift: its curvature is zero and its $q$-summand
is zero. The Tate construction is degreewise split exact on
admissible sequences, proving the last assertions.
\end{proof}

\begin{proposition}[Proved: every odd-primary square-zero torsion direction survives]
\label{C:odd_torsion_allheights}
Let $p$ be odd, let $A$ be a characteristic-zero DVR with maximal
ideal $(p)$, and let $J$ be any $A$-module with $J[p]\ne0$.
For every nonzero $x\in J[p]$, let $L_x$ be the free line over
$k=A\ltimes J$ on which a generator of $C_p$ acts by $1+x$.
Then
\[
 \xi_x=[\operatorname{Ind}_{C_p}^{\mathfrak S_p}L_x]
       -[\operatorname{Ind}_{C_p}^{\mathfrak S_p}\mathbf1]
\]
is nonzero and is outside the degree-$p$ image of $\vartheta_k$.
It vanishes after inverting $p$ and at every residue-field point.
This includes modules $J$ that are $p$-divisible and have torsion.
\end{proposition}
\begin{proof}
First let $K\subseteq J$ be any finitely generated $A$-submodule
containing $x$. The structure theorem over the DVR decomposes
$K$ into a finite free module and cyclic modules $A/p^{r_i}$.
The element $x$ has zero free component, and some cyclic
component is a unit times $p^{r-1}$ in $A/p^r$.
Projection to that component followed by multiplication by a unit
gives an $A$-linear map $\phi:K\to A/p^r$ with
$\phi(x)=p^{r-1}$.
There is a ring homomorphism
\[
 A\ltimes K\longrightarrow A/p^{2r},\qquad
       (a,j)\longmapsto a+p^r\phi(j),
\]
since the image of $K$ is square-zero. It sends $x$ to
$p^{2r-1}$.

Put $m=2r$. Over $A/p^m$ the generator of $L_x$ becomes
$T=1+p^{m-1}$. Lift this scalar to $A$ and take the Tate
differentials $\widetilde d^0=T-1$ and
$\widetilde d^1=1+T+\cdots+T^{p-1}$.
Their common curvature is
\[
 C=\frac{(1+p^{m-1})^p-1}{p^m}\equiv1\pmod p.
\]
The congruence follows by the binomial expansion; for odd $p$
every term after the first is divisible by $p$ after the displayed
division, including the last term when $m=2$.
Thus the $(z-1)$-primary summand is the whole scalar complex, and
\[
 \nu_{z-1,m}(L_x)=m-1\pmod m.
\]
For the twist $g\mapsto g^a$, $1\le a<p$, choose the lift $T^a$.
Its curvature is congruent to $a$ modulo $p$, while
$T^a-1$ still has valuation $m-1$. Exactly the twist $a=1$
contributes to $\nu_{z-1,m}$.

Mackey restriction of the induced representation to $C_p$ gives
each normalizer twist once; every other summand is induced from
the trivial subgroup, so its action is a sum of regular actions
and lifts to $A$. It has zero curvature-length invariant.
Hence $\nu_{z-1,m}$ after this restriction and scalar extension
is $m-1$, a nonzero element of $\mathbb Z/m$.
The ring $A\ltimes K$ is local, so $K_0(A\ltimes K)=\mathbb Z$
and the degree-$p$ polynomial image is generated by Young
permutation classes. All these actions lift to $A$ and have
invariant zero. This proves the assertion for the image of $\xi_x$ over
every such finite submodule $K$.

The ring $A\ltimes J$ is the filtered union of the rings
$A\ltimes K$ over finite submodules $K$ containing $x$.
Proposition~\ref{C:filtered_colimits} applies to the exact
representation groups. If $\xi_x$ were zero, or equal to any
fixed polynomial-image class, that equality would hold at some
such finite stage. The invariant just constructed excludes it.
Finally $x$ is killed by $p$ and is square-zero, so it becomes
zero after inverting $p$ and in every residue field. The claimed
vanishing follows.
\end{proof}

\begin{theorem}[Proved: all square-zero modules over every unramified mixed DVR]
\label{C:allprime_squarezero_criterion}
Let $A$ be any characteristic-zero DVR whose maximal ideal is
$(p)$, for any prime $p$, and let $J$ be any $A$-module. Then
\[
 \begin{aligned}
 \vartheta_{A\ltimes J}\text{ is an isomorphism}
 &\quad\Longleftrightarrow\quad
 \vartheta_{A\ltimes J}\text{ is surjective in degree }p\\
 &\quad\Longleftrightarrow\quad
 p:J\longrightarrow J\text{ is bijective}.
 \end{aligned}
\]
This gives the same intrinsic criterion for every specified
square-zero thickening of $\mathbb Z_{(p)}$, because such a
thickening splits canonically. No finite-generation hypothesis
on $J$ is required.
\end{theorem}
\begin{proof}
For $p=2$ this is Theorem~\ref{C:squarezero_intrinsic_criterion}.
For odd $p$, if multiplication by $p$ is not injective,
Proposition~\ref{C:odd_torsion_allheights} gives failure in degree
$p$. If it is injective, $J$ is torsion-free over the DVR and
Proposition~\ref{C:unramified_flat_squarezero_allprimes} says that
surjectivity is equivalent to $pJ=J$.
When $p$ is bijective on $J$, the nil ideal is uniquely divisible
by every positive integer, so
Proposition~\ref{C:divisible_nil_invariance} and
Proposition~\ref{C:DVR} give the full polynomial answer.
Finally a surjection $k\to\mathbb Z_{(p)}$ with square-zero kernel
has its unique section from $\mathbb Z_{(p)}$: every integer prime
to $p$ is a unit in $k$, because its reduction is a unit and the
kernel is nilpotent. The standard prime-subring map extends
uniquely to that localization and is a section. This identifies
$k$ with the indicated square-zero extension and proves the last
claim.
\end{proof}

The conclusion deliberately separates ramified odd-prime DVRs.
For example a direction killed by a uniformizer $\pi$ when
$p$ has valuation $e>1$ can make the scalar curvature above
divisible by $\pi^{e-1}$. The proof does not assert the same
criterion in that remaining case.
